
\subsubsection{2.1 Controlled Pre-Training Infrastructure}

The Pile \citep{Gao2020} introduced a heterogeneous corpus spanning 22 domains and approximately 800GB of text, including Wikipedia, Books, GitHub, PubMed, StackExchange, and others. It remains the most documented large-scale pre-training corpus in terms of domain composition and proportions.

Pythia \citep{Biderman2023} extended this infrastructure by training 16 autoregressive language models (70M--12B parameters) on The Pile with exact dataloader documentation and releasing 154 intermediate checkpoints per model. The explicit design goal was to enable training dynamics and data attribution studies. Prior work using Pythia has focused on memorization \citep{Biderman2023}, scaling laws, and deduplication effects, not on exploiting the checkpoint trajectory for per-domain exposure analysis. MiniPile \citep{Kaddour2023MiniPile} demonstrated that a 6GB embedding-filtered subset of The Pile retains approximately 98\% of GLUE performance, confirming the importance of data quality filtering, but focused on aggregate performance.

Our approach differs: rather than comparing distinct training runs, we treat the 154-checkpoint trajectory as a panel dataset with naturally occurring domain exposure variation, enabling within-family regression without new training.

\subsubsection{2.2 Domain Mixing Optimization}

DoReMi \citep{Xie2023} uses a small proxy model to find domain weights that minimize excess loss, demonstrating that optimized domain mixing improves aggregate validation loss and downstream benchmarks. RegMix \citep{Liu2024} extends this to a regression-based mixing law, fitting a linear model from domain weights to validation loss at proxy-model scale and extrapolating to larger models. Data Mixing Laws \citep{Ye2024} provide the most direct scaling evidence: domain mixing ratios are explicit predictors in their law formulation.

Our panel regression framework adopts the linear mixing law foundation from RegMix but applies it to individual benchmark scores rather than aggregate validation loss. None of these approaches produce per-domain coefficients for individual benchmarks (MMLU, HellaSwag, ARC, WinoGrande separately). They optimize or characterize aggregate performance. The panel regression framework we implement is the first designed to estimate benchmark-specific domain coefficients from within-family checkpoint variation.

\subsubsection{2.3 Task-Specific Data Selection}

CoLoR-Filter \citep{Brandfonbrener2024} demonstrates that task-conditioned data selection achieves 11--$25\times$ data efficiency for specific downstream benchmarks --- the theoretical motivation underlying domain-benchmark specificity. Topic Over Source \citep{Peng2025} establishes that topic-based mixing consistently outperforms source-based mixing, suggesting that semantic content matters more than source label for training outcomes. This is consistent with the cognitive content proxy approach we employ in h-m1.

\subsubsection{2.4 Benchmark Contamination}

LatestEval \citep{Li2023} and work on evading contamination detection \citep{Dekoninck2024} highlight the difficulty of contamination auditing. We include contamination-adjusted scores via lm-evaluation-harness 13-gram decontamination and document this as a limitation (Section 6.2 L4).

